\chapter{模24 BCD计数器与模60 BCD计数器}

\section{BCD级联的基本思想}

多位BCD计数器 = 多个独立的模10计数器级联。

\begin{center}
\begin{tikzpicture}
  \node[draw, minimum width=2cm, minimum height=1.2cm] (unit) at (0,0) {个位(模10)};
  \node[draw, minimum width=2cm, minimum height=1.2cm] (ten) at (3,0) {十位(模10)};

  \draw[->] (unit.east) -- (ten.west) node[midway,above] {进位};

  \node at (1.5,-1) {\small 个位每计满10，十位加1};
\end{tikzpicture}
\end{center}

级联的本质：低位计数器每完成一轮（回到0），给高位计数器发一个"进位"脉冲，高位计数器加1。

\section{模24 BCD计数器}

模24 BCD = 十位（模2：0-2，计到23后归零）+ 个位（模10：0-9）

等一下——这里有个问题。

\subsection{两种实现方案}

\textbf{方案A：纯BCD级联（个位模10 + 十位模3）}

个位：模10（0-9），满9进位到十位
十位：收到进位+1，但十位只在0-2之间循环（模3）

清零条件：十位=2 且 个位=3（即23），下一拍归零
注意不是十位=2且个位=9（即29）——那会得到30个状态！

\textbf{方案B：两位BCD整体计数，满24清零}

两位合并为8位（十位4位+个位4位），整体+1，到24(0010 0100)时清零。

\subsection{方案A：BCD级联实现}

\begin{lstlisting}[caption={模24 BCD计数器（方案A：级联）}]
library ieee;
use ieee.std_logic_1164.all;
use ieee.std_logic_unsigned.all;

entity bcd_counter_mod24 is
  port (
    clk   : in  std_logic;
    rst_n : in  std_logic;
    tens  : out std_logic_vector(3 downto 0);
    units : out std_logic_vector(3 downto 0)
  );
end bcd_counter_mod24;

architecture behavioral of bcd_counter_mod24 is
  signal units_cnt : std_logic_vector(3 downto 0);
  signal tens_cnt  : std_logic_vector(3 downto 0);
  signal units_carry : std_logic;
begin
  process(clk, rst_n)
  begin
    if rst_n = '0' then
      units_cnt <= (others => '0');
      tens_cnt  <= (others => '0');
      units_carry <= '0';

    elsif rising_edge(clk) then
      -- 检测满24清零
      if tens_cnt = "0010" and units_cnt = "0011" then  -- 23
        units_cnt <= (others => '0');
        tens_cnt  <= (others => '0');
        units_carry <= '0';
      else
        -- 个位计数
        if units_cnt = "1001" then  -- 9
          units_cnt <= (others => '0');
          units_carry <= '1';
        else
          units_cnt <= units_cnt + 1;
          units_carry <= '0';
        end if;

        -- 十位计数（收到进位时）
        if units_carry = '1' then
          tens_cnt <= tens_cnt + 1;
        end if;
      end if;
    end if;
  end process;

  tens  <= tens_cnt;
  units <= units_cnt;
end behavioral;
\end{lstlisting}

\subsection{方案B：统一计数}

\begin{lstlisting}[caption={模24 BCD计数器（方案B：统一进位判断）}]
-- 更简洁的实现：直接在同一个process中处理
process(clk, rst_n)
begin
  if rst_n = '0' then
    units_cnt <= (others => '0');
    tens_cnt  <= (others => '0');
  elsif rising_edge(clk) then
    -- 先判断是否清零
    if tens_cnt = "0010" and units_cnt = "0011" then  -- 23
      units_cnt <= (others => '0');
      tens_cnt  <= (others => '0');
    -- 个位进位判断
    elsif units_cnt = "1001" then  -- 9
      units_cnt <= (others => '0');
      tens_cnt  <= tens_cnt + 1;
    else
      units_cnt <= units_cnt + 1;
    end if;
  end if;
end process;
\end{lstlisting}

\section{模60 BCD计数器}

模60 BCD = 十位（模6：0-5）+ 个位（模10：0-9）

清零条件：十位=5 且 个位=9（即59）

\begin{lstlisting}[caption={模60 BCD计数器}]
library ieee;
use ieee.std_logic_1164.all;
use ieee.std_logic_unsigned.all;

entity bcd_counter_mod60 is
  port (
    clk   : in  std_logic;
    rst_n : in  std_logic;
    tens  : out std_logic_vector(3 downto 0);
    units : out std_logic_vector(3 downto 0)
  );
end bcd_counter_mod60;

architecture behavioral of bcd_counter_mod60 is
  signal units_cnt : std_logic_vector(3 downto 0);
  signal tens_cnt  : std_logic_vector(3 downto 0);
begin
  process(clk, rst_n)
  begin
    if rst_n = '0' then
      units_cnt <= (others => '0');
      tens_cnt  <= (others => '0');
    elsif rising_edge(clk) then
      -- 满59清零
      if tens_cnt = "0101" and units_cnt = "1001" then  -- 59
        units_cnt <= (others => '0');
        tens_cnt  <= (others => '0');
      -- 个位满9进位
      elsif units_cnt = "1001" then
        units_cnt <= (others => '0');
        tens_cnt  <= tens_cnt + 1;
      else
        units_cnt <= units_cnt + 1;
      end if;
    end if;
  end process;

  tens  <= tens_cnt;
  units <= units_cnt;
end behavioral;
\end{lstlisting}

\section{BCD级联的通用模板}

\begin{designflow}[多位BCD计数器通用模板]
\textbf{个位计数器}：
\begin{itemize}
  \item 模10：计到9后清零，产生进位
\end{itemize}

\textbf{十位计数器}：
\begin{itemize}
  \item 收到进位后+1
  \item 有自己独立的模值限制（取决于应用）
  \item 模24：十位模3（0-2），满23清零
  \item 模60：十位模6（0-5），满59清零
\end{itemize}

\textbf{整体清零条件}（同时清除个位和十位）：
\begin{itemize}
  \item 模24：十位=2 且 个位=3
  \item 模60：十位=5 且 个位=9
  \item 模12：十位=1 且 个位=1（如果要做成12的BCD）
\end{itemize}
\end{designflow}

\section{数字钟表的完整结构}

一个数字钟表 = 三个BCD计数器级联：

\begin{center}
\begin{tikzpicture}
  \node[draw, minimum width=2.5cm, minimum height=1.2cm] (sec) at (0,0) {模60秒};
  \node[draw, minimum width=2.5cm, minimum height=1.2cm] (min) at (3.5,0) {模60分};
  \node[draw, minimum width=2.5cm, minimum height=1.2cm] (hour) at (7,0) {模24时};

  \draw[->] (sec.east) -- (min.west) node[midway,above] {进位};
  \draw[->] (min.east) -- (hour.west) node[midway,above] {进位};

  \node at (3.5,-1.5) {\small 秒满60→分+1, 分满60→时+1, 时满24→清零};
\end{tikzpicture}
\end{center}

这正是模24和模60 BCD计数器的典型应用场景。

\section{计数器部分总结}

\begin{keypoint}[计数器统一框架总结]
\textbf{所有计数器本质上只有一个公式：}

在时钟上升沿：$cnt \leftarrow cnt + 1$（如果$cnt = N-1$则$cnt \leftarrow 0$）

\begin{itemize}
  \item 二进制计数器：一个整体二进制数（4/5/6位）
  \item BCD计数器：每个十进制位独立（个位+十位各4位）
  \item 区别：表示方式不同，但核心转移逻辑完全相同
  \item 模$2^N$计数器：唯一不需要显式清零的（自然溢出）
  \item 所有其他模值：都需要"if cnt=N-1 then cnt<=0"的显式清零
\end{itemize}
\end{keypoint}
